\documentclass[11pt]{article}
\usepackage{graphicx}
\usepackage{float}
\usepackage[margin=1in]{geometry}
\usepackage{amsmath, amssymb}
\usepackage{booktabs}

\begin{document}

\section*{setup}
\begin{itemize}
\item Chain $n=25$, Case~0, $\gamma=0.1$, $dt=0.001$. NESS: $T_1=10$, $T_{25}=2$;
equilibrium: $T_1=T_{25}=6$. SIMD runs, $1.01\times 10^9$ samples per site pair,
site pairs $(I_j, I_{j+1}, \theta_j)$ at $j=6,12,18$.
\item Histograms: $80^3$ bins on $[0,4]^2\times[-\pi,\pi)$. All .hist
files now use the dense format (with all $512{,}000$ bins , zeros
included).
\item Notation: $p_\beta$ = equilibrium marginal at temperature
$\beta$. $q$ = NESS marginal. 

\end{itemize}

\noindent \textbf{1. Compare marginal PDF of $(I_a, I_{a+1}, \theta_a)$ and the 2-mode Gibbs by comparing $-\log(PDF)$ and $\frac{1}{2} (I_1 + I_2)^2 - \frac{1}{4} (I_1^2 + I_2^2) + I_1I_2\cos \theta_1$. 
}

\begin{itemize}
    \item $-\log p$ v.s. $H_{\text{2-mode}} = \tfrac12 (I_a+I_b)^2 - \tfrac14 (I_a^2+I_b^2)
+ I_a I_b\cos\theta$
at fixed $\theta$ slices ($\theta = 0, \pi/2, \pi$, $\cos\theta = -1, 0, 1$) 

\end{itemize}
\begin{figure}[H]
    \centering
    \includegraphics[width=0.8\linewidth]{compare_eq.png}
\end{figure}


\begin{itemize}
    \item $-\log q$ v.s. $H_{\text{2-mode}} = \tfrac12 (I_a+I_b)^2 - \tfrac14 (I_a^2+I_b^2)
+ I_a I_b\cos\theta$
at fixed $\theta$ slices ($\theta = 0, \pi/2, \pi$, $\cos\theta = -1, 0, 1$) 


\end{itemize}

\begin{figure}[H]
    \centering
    \includegraphics[width= 1\linewidth]{6_5/compare_neq.png}
\end{figure}



\noindent \textbf{2. Compare marginal PDF of $(I_a, I_{a+1}, \theta_a)$ of NESS and the same marginal PDF of equilibrium by comparing the $-\log (PDF)$ of both. }

\begin{itemize}
    \item \textbf{Equilibrium: $p_\beta(I_a, I_{a+1}, \theta_a) : = \int_A \exp(-H/\beta) \mathrm{d}x$, where $A$ is the collection of variables other than $(I_a, I_{a+1}, \theta_a)$.}
    \item \textbf{Nonequilibrium: Does the marginal PDF of $(I_a, I_{a+1}, \theta_a)$, denoted by $q(I_a, I_{a+1}, \theta_a)$ approach to $p_{\beta_0}(I_a, I_{a+1}, \theta_a)$ of some $\beta_0$ as the length of the chain ($n$) increases? How does $\beta_0$ depend on $(T_L, T_R, a/n)$? }
    \item \textbf{Test: fix $a/n$. Let $n$ increase. Let $p_\beta$ be the equilibrium with boundary temperature $\beta$. Find the minimum of $\min_x\|\exp(x\log p) - q\|$. So the effective boundary temperature is $\beta_0:= \beta/x$. Compare $p_{\beta_0}$ and $q$.}
\end{itemize}






\paragraph{(i) Find $\min_x \|\exp(x\log p_6) - q\|$ and the effective temperature
$\beta_0 = 6/x$.}



We fit $\log q = x\,\log p_6 + c$ over the common support (count $>50$ in
both histograms). Results at fixed $a/n$:
\begin{center}
\begin{tabular}{lcccccc}
\toprule
 & \multicolumn{3}{c}{$n=25$ ($a=6,12,18$)} & \multicolumn{3}{c}{$n=50$ ($a=12,24,36$)} \\
 & $x$ & $R^2$ & $T_0$ & $x$ & $R^2$ & $T_0$ \\
\midrule
$a/n=0.24$ & $0.7800$ & $0.9914$ & $7.69$ & $0.8408$ & $0.9929$ & $7.14$ \\
$a/n=0.48$ & $0.8432$ & $0.9763$ & $7.12$ & $0.8634$ & $0.9908$ & $6.95$ \\
$a/n=0.72$ & $0.9021$ & $0.9314$ & $6.65$ & $0.8667$ & $0.9453$ & $6.92$ \\
\bottomrule
\end{tabular}
\end{center}

(next page)

\newpage
\paragraph{(ii) Compare $p_{\beta_0}$ and $q$.}


For each $(n, a)$, compare the rescaled equilibrium $p_6^{\,x}$ with $q$ at fixed $\theta$ slices.


\begin{figure}[H]
    \centering
    \includegraphics[width=1\linewidth]{n25_a12.png}
    \includegraphics[width=1\linewidth]{n25_a18.png}
    \includegraphics[width=0.8\linewidth]{n50_a24.png}
\end{figure}

\textbf{Control: }equilibrium simulated at $\beta = 7.12$
($= T_0(n{=}25, a{=}12)$).
Fitting $\log p_{7.12} = x \log p_6 + c$ gives $x = 0.9220$ at all three
sites (spread $0.03\%$), $R^2 = 0.9975$. This differs from
$6/7.12 = 0.843$. Fitting $q(a{=}12)$ against $p_{7.12}$ gives $x = 0.905$, not $1$. Guess: close to
$\sqrt{6/7.12} = 0.918$, might come from the mean field
$\langle \sum_j I_j \rangle \propto \sqrt{nT}$ changing with temperature in the full equilibrium? 


\paragraph{(iii) What $\beta_0$ is: dependence on $(T_L, T_R, a/n)$.}

\begin{itemize}
    \item \textbf{Dependence on $a/n$}

We estimate the local temperature of site $a$
from the $\langle I_j \rangle$ profiles:
\[
T_{\mathrm{kin}}(a) = 6 \cdot
\frac{\langle I_a\rangle_{\mathrm{NESS}}}{\langle I_a\rangle_{\mathrm{Eq}}},
\]
averaged over the pair $(a, a{+}1)$.
Comparing it with $T_0^{\mathrm{box}}$ from the fit: 


$T_0^{\mathrm{box}}$: same fit, count-weighted, restricted to $I_a, I_b < 2.5$ (weighting the core of the distribution rather than the tail).

\begin{center}
\begin{tabular}{lcccccc}
\toprule
$a/n$ & \multicolumn{2}{c}{$n=25$} & \multicolumn{2}{c}{$n=50$} & \multicolumn{2}{c}{$n=100$} \\
 & $T_{0}^{\mathrm{box}}$ & $T_{\mathrm{kin}}$ & $T_{0}^{\mathrm{box}}$ & $T_{\mathrm{kin}}$ & $T_{0}^{\mathrm{box}}$ & $T_{\mathrm{kin}}$ \\
\midrule
$0.24$ & $7.32$ & $7.68$ & $7.26$ & $7.79$ & $7.12$ & $7.81$ \\
$0.48$ & $6.41$ & $6.85$ & $6.45$ & $7.04$ & $6.45$ & $6.99$ \\
$0.72$ & $5.49$ & $5.80$ & $5.66$ & $6.05$ & $5.76$ & $6.02$ \\
\bottomrule
\end{tabular}
\end{center}
  
\end{itemize}


$T_0$ follows the local kinetic temperature at both chain lengths
(it is below $T_{\mathrm{kin}}$ by $5$--$7\%$ at every site). So, to leading
order, the dependence of $\beta_0$ on $(T_L, T_R, a/n)$ is simply the shape
of the temperature profile.


\begin{itemize}
    \item \textbf{Dependence on $(T_L, T_R)$}

Runs at $(8,4)$, $(9,3)$, $(10,2)$
(same mean temperature $6$, same $p_6$ baseline), $n=25$:
\begin{center}
\begin{tabular}{llccc}
\toprule
$(T_L,T_R)$ & $a$ & $x$ & $R^2$ & $T_0$ \\
\midrule
$(8,4)$  & $6$  & $0.8785$ & $0.9950$ & $6.83$ \\
         & $12$ & $0.9366$ & $0.9906$ & $6.41$ \\
         & $18$ & $1.0008$ & $0.9777$ & $6.00$ \\
$(9,3)$  & $6$  & $0.8230$ & $0.9928$ & $7.29$ \\
         & $12$ & $0.8872$ & $0.9819$ & $6.76$ \\
         & $18$ & $0.9498$ & $0.9512$ & $6.32$ \\
$(10,2)$ & $6$  & $0.7800$ & $0.9914$ & $7.69$ \\
         & $12$ & $0.8432$ & $0.9763$ & $7.12$ \\
         & $18$ & $0.9021$ & $0.9314$ & $6.65$ \\
\bottomrule
\end{tabular}
\end{center}
    
\end{itemize}





\paragraph{(iv) Does $q \to p_{\beta_0}$ as $n$ increases (fixed $a/n$)?}

From $n=25$ to $n=50$, the fit improves at every site:
\begin{center}
\begin{tabular}{lcc}
\toprule
$a/n$ & $R^2$ ($n=25$) & $R^2$ ($n=50$) \\
\midrule
$0.24$ & $0.9914$ & $0.9929$ \\
$0.48$ & $0.9763$ & $0.9908$ \\
$0.72$ & $0.9314$ & $0.9453$ \\
\bottomrule
\end{tabular}
\end{center}

$R^2$ increases at every $a/n$. For reference, the comparison of the two equilibria ($p_{7.12}$ vs.\ $p_6$, $R^2 = 0.9975$)
sets the ceiling for this fit. at $a/n = 0.24$ the NESS fit at $n=50$ is
already close to this level. At $n=100$ the core (box) fit is consistent with this trend (see (iii)); the full-support $R^2$ is tail-limited at $dt=0.001$ (see Note).

\paragraph{Note ($n=100$)}
The three sites at $n=100$ give nearly the same $x$ ($T_0 \approx 6.4$) while the profiles show a clear gradient. 

1. rerun $n=100$ with a
corrected initialization (chains start at the stationary action scale,
verified by monitoring $\sum_j I_j$ during burn-in) and the fits are still the same. 2. Change the time
step ($dt: 0.001 \to 0.0005$) roughly halves the histogram overflow rate
($0.13\% \to 0.06\%$), and the same tail show up in the equilibrium
run, so it might be the Euler-integration
artifact at $dt = 0.001$, not physics.

The artifact is buried in the wide thermal tails at $n=25$
and $n=50$, so those results are unaffected. 

\paragraph{Code}
\begin{itemize}
\item \texttt{lte\_histogram\_simd.cpp}: SIMD simulation; produces all
histogram files (\texttt{*.hist}) and the $\langle I_j\rangle$ profile files.
\item \texttt{compare.m}: Q1 figures ($-\log$ PDF vs.\ $H_{\text{2-mode}}$
at fixed $\theta$ slices).
\item \texttt{fit.m}: all Q2 fits and figures, tables in
(i), (iii), (iv) ($x$, $R^2$, $T_0$, $T_0^{\mathrm{box}}$), the control fits
in (ii), and the heatmap figures. $T_{\mathrm{kin}}$ in (iii) is computed
directly from the profile files.
\end{itemize}



\end{document}